\chapter*{پیش‌گفتار: کلان‌داده و مبانی محاسبات مقیاس‌پذیر}
\addcontentsline{toc}{chapter}{پیش‌گفتار: کلان‌داده و مبانی محاسبات مقیاس‌پذیر}
\markboth{پیش‌گفتار}{پیش‌گفتار}

\section*{مقدمه و انگیزه}
در دهه‌های اخیر، با رشد تصاعدی وب، شبکه‌های اجتماعی، حسگرهای اینترنت اشیاء\footnote{\lr{Internet of Things (IoT)}}، معاملات مالی و سامانه‌های پایش علمی، جوامع بشری وارد عصری شده‌اند که حجم و سرعت تولید داده‌ها از توان تحلیلی و ذخیره‌سازی ابزارهای سنتی پایگاه داده فراتر رفته است. این پدیده تحت عنوان \textbf{کلان‌داده}\footnote{\lr{Big Data}} شناخته می‌شود. 

مواجهه با کلان‌داده صرفاً یک مسئله ارتقای سخت‌افزاری نیست؛ بلکه نیازمند بازنگری بنیادین در نظریه الگوریتم‌ها، مدل‌های ریاضی و معماری سیستم‌های محاسباتی است. الگوریتم‌هایی که دارای مرتبه زمانی درجه دوم $(O(n^2))$ یا حتی درجه اول خطی همراه با دسترسی تصادفی به حافظه هستند، بر روی کلان‌داده‌هایی با مقیاس ترابایت و پتابایت عملاً ناکارآمد بوده و متوقف می‌شوند. از این رو، چارچوب نظری جدیدی لازم است که الگوریتم‌های زیرخطی\footnote{\lr{Sublinear Algorithms}}، الگوریتم‌های برخط جریانی\footnote{\lr{Streaming Algorithms}}، روش‌های تصادفی‌سازی و تقریب\footnote{\lr{Randomized and Approximation Algorithms}}، و پردازش موازی و توزیع‌شده\footnote{\lr{Distributed and Parallel Computing}} را در کانون توجه قرار دهد.

\section*{مرجع اصلی و رویکرد آموزشی این اثر}
این کتاب و مجموعه یادداشت‌های جامع درسی بر پایه سرفصل‌های رسمی دوره کارشناسی ارشد مهندسی کامپیوتر و علم داده و با تکیه بر مرجع پیشگام و معتبر دانشگاه استنفورد، یعنی کتاب:
\begin{center}
\textbf{Mining of Massive Datasets \enparen{MMDS}}~\cite{leskovec2020mining} \\
\textit{Jure Leskovec, Anand Rajaraman, Jeffrey D. Ullman}
\end{center}
تدوین شده است. در تدوین این مجموعه، اصول آموزشی زیر به عنوان راهنمای عمل مد نظر قرار گرفته‌اند:

\begin{enumerate}
    \item \textbf{محوریت اسلایدهای درسی در ساختار و جهت‌گیری}: ساختار هر فصل، تقدم و تاخر مباحث و میزان تاکید بر مفاهیم دقیقاً منطبق بر اسلایدهای رسمی درس تنظیم شده است تا دانشجویان پیوستگی کاملی میان جلسات درس و متن مکتوب احساس کنند.
    \item \textbf{جامعیت و خودبسندگی با بهره‌گیری از متن مرجع}: اسلایدهای درسی به دلیل ماهیت ارائه شفاهی، غالباً موجز و حاوی نکات کلیدی هستند. در این اثر، با رجوع دقیق به متن کتاب مرجع، کلیه تعاریف ریاضی، اثبات‌های راهگشا، تحلیل پیچیدگی الگوریتم‌ها و مثال‌های عددی گام‌به‌گام تشریح شده‌اند به گونه‌ای که مطالعه این کتاب بی‌نیاز از بازخوانی مدام اسلایدها باشد.
    \item \textbf{نگارش دانشگاهی و اصطلاح‌شناسی دقیق}: متن با نگارش علمی، روان و استاندارد فارسی نگاشته شده و برای تمامی اصطلاحات تخصصی، در اولین کاربرد، معادل انگلیسی آن‌ها در پاورقی ذکر شده است تا هم غنای علمی زبان فارسی حفظ شود و هم ارتباط دانشجو با مقالات و ادبیات بین‌المللی گسسته نگردد.
    \item \textbf{تمرکز بر بینش شهودی و محاسبات واقعی}: ریاضیات و فرمول‌ها هرگز بدون توضیح شهودی رها نشده‌اند؛ برای هر رابطه ریاضی، معنای هر متغیر، دلیل بهره‌گیری از فرمول و چگونگی تفسیر نتایج تشریح گردیده است.
\end{enumerate}

\section*{نقشه راه فصول کتاب}
محتوای این اثر در سیزده فصل اصلی و بر اساس یک جریان منطقی منسجم به شرح زیر سازمان‌دهی شده است:

\begin{itemize}
    \item \textbf{فصل ۱: داده‌کاوی و فضاهای با بعد بالا} -- بررسی ماهیت داده‌کاوی، مقایسه رویکرد یادگیری ماشین در برابر مدل‌های آماری، اصل بونفرونی\footnote{\lr{Bonferroni's Principle}} و اجتناب از کشف الگوهای کاذب، و چالش‌های محاسباتی در فضاهای ابعادی بزرگ.
    \item \textbf{فصل ۲: نگاشت-کاهش و پشته نرم‌افزاری} -- معماری سیستم‌های فایل توزیع‌شده (مانند \lr{HDFS})، مدل برنامه‌نویسی نگاشت-کاهش\footnote{\lr{MapReduce}}، شیوه‌های مدیریت خطا، و پیاده‌سازی عملیات ماتریسی و جبر رابطه‌ای در مقیاس وسیع.
    \item \textbf{فصل ۳: یافتن آیتم‌های مشابه} -- مسئله مقایسه اسناد در وب، روش شینگلینگ\footnote{\lr{Shingling}}، معیار تشابه جاکارد\footnote{\lr{Jaccard Similarity}}، تکنیک \lr{Min-Hashing}، و درهم‌سازی حساس به مکان\footnote{\lr{Locality-Sensitive Hashing (LSH)}} جهت کاهش ابعاد جستجو.
    \item \textbf{فصل ۴: کاوش جریان‌های داده} -- مدل محاسباتی جریان داده، نمونه‌گیری مخزنی\footnote{\lr{Reservoir Sampling}}، فیلترهای بلوم\footnote{\lr{Bloom Filters}}، تخمین تعداد اعضای متمایز با الگوریتم فلژوله-مارتین\footnote{\lr{Flajolet-Martin Algorithm}}، و شمارش بیت‌های یک در پنجره لغزان با الگوریتم \lr{DGIM}.
    \item \textbf{فصل ۵: تحلیل پیوندها و رتبه‌بندی وب} -- ساختار گراف وب، الگوریتم مشهور رتبه‌بندی صفحه\footnote{\lr{PageRank}}، تله‌های عنکبوتی و بن‌بست‌ها، راهکار بازآغازی تصادفی، رتبه‌بندی وابسته به موضوع، و الگوریتم \lr{HITS}.
    \item \textbf{فصل ۶: مجموعه‌آیتم‌های پرتکرار} -- مدل سبد بازار\footnote{\lr{Market-Basket Model}}، قوانین انجمنی\footnote{\lr{Association Rules}}، الگوریتم \lr{A-Priori} و بهینه‌سازی‌های حافظه‌محور نظیر \lr{PCY} و \lr{Multistage}، و الگوریتم‌های توزیع‌شده \lr{SON} و \lr{Toivonen}.
    \item \textbf{فصل ۷: خوشه‌بندی در ابعاد بالا} -- محدودیت‌های روش‌های خوشه‌بندی سنتی در ابعاد بالا، الگوریتم‌های سلسله‌مراتبی، \lr{K-Means} و بهینه‌سازی‌های \lr{BFR} و \lr{CURE} برای خوشه‌بندی داده‌های عظیم با شکل‌های پیچیده.
    \item \textbf{فصل ۸: تبلیغات برخط در وب} -- مدل تطابق دوبخشی برخط\footnote{\lr{Online Bipartite Matching}}، نسبت رقابتی\footnote{\lr{Competitive Ratio}}، الگوریتم \lr{BALANCE} و سازوکار حراج آگهی گوگل \enparen{AdWords}.
    \item \textbf{فصل ۹: سامانه‌های توصیه‌گر} -- ماتریس سودمندی\footnote{\lr{Utility Matrix}}، فیلترینگ محتوامحور، فیلترینگ مشارکتی\footnote{\lr{Collaborative Filtering}}، و تجزیه ماتریسی با کاهش مقدار منفرد \enparen{SVD} و فاکتورسازی ماتریس \lr{UV}.
    \item \textbf{فصل ۱۰: کاوش گراف‌های شبکه‌های اجتماعی} -- ساختار جوامع در شبکه‌های اجتماعی، مرکزیت بینابینی و الگوریتم گیروان-نیومن\footnote{\lr{Girvan-Newman Algorithm}}، معیار پیمانه‌ای\footnote{\lr{Modularity}}، و خوشه‌بندی طیفی\footnote{\lr{Spectral Clustering}} با ماتریس لاپلاسین.
    \item \textbf{فصل ۱۱: کاهش بعد} -- تجزیه مقادیر منفرد\footnote{\lr{Singular Value Decomposition (SVD)}}، تحلیل مؤلفه‌های اصلی\footnote{\lr{Principal Component Analysis (PCA)}}، لم جانسون-لیندن‌اشتراوس\footnote{\lr{Johnson-Lindenstrauss Lemma}}، و تجزیه ماتریسی \lr{CUR}.
    \item \textbf{فصل ۱۲: یادگیری ماشین در مقیاس بزرگ} -- ماشین‌های بردار پشتیبان\footnote{\lr{Support Vector Machines (SVM)}}، تابع هزینه لولا\footnote{\lr{Hinge Loss}}، گرادیان کاهشی تصادفی\footnote{\lr{Stochastic Gradient Descent (SGD)}} با یک گذر، و موازی‌سازی الگوریتم‌های یادگیری.
    \item \textbf{فصل ۱۳: شبکه‌های عصبی و یادگیری عمیق} -- پرسپترون چندلایه، پس‌انتشار خطا\footnote{\lr{Backpropagation}}، شبکه‌های پیچشی \enparen{CNN} و ساختارهای توجه و ترنسفورمرها در کلان‌داده و پردازش‌های توزیع‌شده \lr{GPU}.
\end{itemize}

\begin{summarybox}[راهنمای مطالعه و آمادگی برای آزمون]
برای بهره‌گیری حداکثری از این کتاب توصیه می‌شود:
\begin{enumerate}
    \item در ابتدای هر فصل، بخش «مقدمه و اهداف فصل» را مرور فرمایید تا نقشه مفهومی مبحث در ذهن شکل گیرد.
    \item به کادرهای مشخص‌شده با رنگ‌های مختلف توجه فرمایید:
    \begin{itemize}
        \item کادرهای \textbf{آبی} بیانگر تعاریف پایه و اصطلاحات رسمی هستند.
        \item کادرهای \textbf{سبز} قضایا، لم‌ها و گزاره‌های ریاضی اثبات‌شده هستند.
        \item کادرهای \textbf{خاکستری} مشخصات گام‌به‌گام و شبه‌کد الگوریتم‌ها هستند.
        \item کادرهای \textbf{نارنجی} مثال‌های عددی حل‌شده را در بر دارند.
        \item کادرهای \textbf{قرمز} حاوی نکات طلایی امتحانی و دام‌های متداول مفهومی می‌باشند.
    \end{itemize}
    \item مثال‌های عددی را به صورت دستی روی کاغذ دنبال کنید تا مهارت لازم در آزمون‌ها و تمرین‌های عملی حاصل گردد.
\end{enumerate}
\end{summarybox}

\begin{summarybox}[شناسنامه علمی و عوامل تدوین اثر]
\begin{itemize}
    \item \textbf{عنوان کتاب:} مبانی و الگوریتم‌های کلان‌داده \enparen{Mining of Massive Datasets}
    \item \textbf{گردآورنده و پژوهشگر:} \textbf{احسان شهبازی} \enparen{Ehsan Shahbazi}
    \item \textbf{پست الکترونیک:} \href{mailto:ehsan.shahbazipc@gmail.com}{\lr{ehsan.shahbazipc@gmail.com}}
    \item \textbf{مخزن و صفحه رسمی گیت‌هاب:} \href{https://github.com/EhsanShahbazii}{\lr{github.com/EhsanShahbazii}}
    \item \textbf{مرجع درسی:} کتاب \lr{Mining of Massive Datasets}، دانشگاه استنفورد
    \item \textbf{مقطع و جامعه مخاطب:} دانشجویان تحصیلات تکمیلی مهندسی کامپیوتر، هوش مصنوعی و علم داده
\end{itemize}
\end{summarybox}
